In the last decades, generic modules and generic bricks have caught the interest of researchers thanks to their connections with the representation-type and properties of finite-dimensional algebras. We recall that a module over a finite-dimensional algebra is called \emph{generic} if it is indecomposable with infinite length but finite endolength, that is its length as a module over its endomorphism ring. If the generic module is also a \emph{brick}, that is its endomorphism ring is a division ring, then it is called \emph{generic brick}.

Crawley-Boevey showed that representation-infinite algebras are characterised by the existence of a generic module \cite[Theorem 4.5]{crawley1991tame} and that tame algebras are equivalent to generically tame algebras \cite[Theorem 4.4]{crawley1991tame}.
It was conjectured by Mousavand and Paquette \cite[Conjecture 4.1]{mousavand2025biserial} that for any algebra, being \emph{brick-infinite}, being \emph{brick-continuous} or admitting a generic brick were three equivalent properties (where an algebra is brick-infinite if it admits an infinite family of non-isomorphic bricks, and brick-continuous if it admits an infinite family of non-isomorphic bricks with the same dimension). They proved such result holds for biserial algebras \cite[Corollary 4.5]{mousavand2025biserial}.
Then, Bautista, Pérez and Salmerón proved the equivalence of brick-continuity and existence of a generic brick for tame algebras \cite[Corollary 1.3]{bau:gen}. After that, Schlegel gave an alternative proof of such result \cite[Theorem 5.6]{schlegel2026infinitetautiltingtheory} and proved the full conjecture for algebras whose Krull-Gabriel dimension is defined \cite[Theorem 5.8]{schlegel2026infinitetautiltingtheory}.\medskip


In this paper we will discuss the existence of brick-continuous families and generic bricks for representation-infinite incidence algebras of posets. The \emph{incidence algebra} $\incidencealgebra{P}$ of a finite poset $P$ will be defined as the bound path $\mathbf{k}$-algebra $\mathbf{k}Q_P/I_P$, where $Q_P$ is the \emph{Hasse quiver} of $P$ and $I_P$ the ideal of $\mathbf{k}Q_P$ generated by the relations $\{p_1 - p_2 \}_{(p_1,p_2)}$, where we index over all pairs of paths with the same start and endpoints.

After giving the necessary preliminary notions, we will focus on the twelve classes of \emph{minimal representation-infinite incidence algebras} given by the so called \emph{critical} posets \cite{lou:ind, lou:repr}.
Since these twelve families are representation-infinite they admit a generic module by \cite[Theorem 4.5]{crawley1991tame}. Moreover, such algebras are $\tau$-tilting infinite, thanks to \cite[Theorem 2.15]{borve:tau}. Hence, by \cite[Theorem 4.14]{sent:brick}, they admit a brick that is not finitely generated. We can therefore ask whether such algebras admit a module having both properties, that is a generic brick.

We will show that each class has tame representation type and it is brick-continuous, thus they admit a generic brick (see \cite[Corollary 1.3]{bau:gen} and \cite[Theorem 5.6]{schlegel2026infinitetautiltingtheory}). This process will utilise the concept of Tits quadratic form together with its \emph{radical vectors}. We will give a method to explicitly construct the brick-continuous families and a generic brick for each of the twelve classes and thanks to the properties of APR-tilting functors, we will prove the following result:

\begin{reptheorem}{thm:genbrickcritical}
	\emph{The twelve families of minimally representation-infinite incidence $\K$-algebras given by critical posets in Table \ref{tab:Loupias frames} (considering any orientation of the undirected edges), and also the opposites of these, are brick-continuous and admit a generic brick which can be explicitly constructed.}
\end{reptheorem}

After that, we will further expand our view. Indeed, any representation-infinite poset can be reduced to one of the twelve classes (see Lemma \ref{lemma:infinitereduction}) through the reduction techniques devised by Loupias \cite{lou:ind,lou:repr}. We will show that this reduction process can be ``reversed'' in order to preserve brick-continuity and generic bricks using what we will call \emph{extension functors}. This will lead to our second main result:

\begin{reptheorem}{thm:genbrickinfinite}
	\emph{Any representation-infinite poset $P$ is brick-continuous and admits a generic brick over $\incidencealgebra{P}$ that can be explicitly computed.}
\end{reptheorem}

Lastly, using analogous ideas, we will generalise this result to a wider class of algebras:

\begin{reptheorem}{thm:generalisation}
	\emph{Let $\Lambda, \Lambda'$ be finite-dimensional algebras.
	\begin{enumerate}
		\item If $\Lambda'$ is brick-continuous and $F \colon \Modf{\Lambda'} \to \Modf{\Lambda}$ is a fully faithful exact functor, then $\Lambda$ is brick-continuous.
		\item If $\Lambda'$ has a generic brick and $F \colon \Mod{\Lambda'} \to \Mod{\Lambda}$ is a fully faithful exact functor that preserves endofinite modules, then $\Lambda$ admits a generic brick.
	\end{enumerate}
	}
\end{reptheorem}

A direct application of this result is given when $\Lambda'=\incidencealgebra{P}$ for a representation-infinite poset $P$, for example having a ring epimorphism $f\colon \Lambda \to \incidencealgebra{P}$.
